\documentclass[aspectratio=169]{beamer}
\usepackage[english]{babel}
\usepackage[utf8]{inputenc}
\usepackage[T1]{fontenc}
\usepackage{lmodern}
\usepackage{listings}
\usepackage{xcolor}
\usepackage{graphicx}
\usepackage{textcomp}
\DeclareUnicodeCharacter{2014}{---}
\DeclareUnicodeCharacter{2013}{--}
\DeclareUnicodeCharacter{2212}{-}
\DeclareUnicodeCharacter{2192}{\ensuremath{\rightarrow}}
\DeclareUnicodeCharacter{00B1}{\ensuremath{\pm}}
\DeclareUnicodeCharacter{00D7}{\ensuremath{\times}}
\DeclareUnicodeCharacter{00B7}{\ensuremath{\cdot}}
\DeclareUnicodeCharacter{20AC}{EUR}
\DeclareUnicodeCharacter{2070}{\textsuperscript{0}}
\DeclareUnicodeCharacter{00B9}{\textsuperscript{1}}
\DeclareUnicodeCharacter{00B2}{\textsuperscript{2}}
\DeclareUnicodeCharacter{00B3}{\textsuperscript{3}}
\DeclareUnicodeCharacter{2074}{\textsuperscript{4}}
\DeclareUnicodeCharacter{2075}{\textsuperscript{5}}
\DeclareUnicodeCharacter{2076}{\textsuperscript{6}}
\DeclareUnicodeCharacter{2077}{\textsuperscript{7}}
\DeclareUnicodeCharacter{2078}{\textsuperscript{8}}
\DeclareUnicodeCharacter{2079}{\textsuperscript{9}}
\DeclareUnicodeCharacter{2248}{\ensuremath{\approx}}
\DeclareUnicodeCharacter{2264}{\ensuremath{\le}}

% Colors for code
\definecolor{codebg}{RGB}{245,245,245}
\definecolor{codecomment}{RGB}{0,128,0}
\definecolor{codekeyword}{RGB}{0,0,255}
\definecolor{codestring}{RGB}{163,21,21}
\definecolor{bandcolor}{RGB}{200,50,50}

\lstdefinestyle{pythonstyle}{
    language=Python,
    backgroundcolor=\color{codebg},
    basicstyle=\ttfamily\scriptsize,
    keywordstyle=\color{codekeyword},
    commentstyle=\color{codecomment},
    stringstyle=\color{codestring},
    showstringspaces=false,
    breaklines=true,
    frame=single,
    rulecolor=\color{black},
    escapeinside={(*@}{@*)},
    moredelim=**[l][\color{bandcolor}\bfseries]{\#\ ---\ NEW},
    moredelim=**[l][\color{bandcolor}\bfseries]{\#\ ---\ CHANGED},
    literate=
      {á}{{\'a}}1 {é}{{\'e}}1 {í}{{\'i}}1 {ó}{{\'o}}1 {ú}{{\'u}}1
      {Á}{{\'A}}1 {É}{{\'E}}1 {Í}{{\'I}}1 {Ó}{{\'O}}1 {Ú}{{\'U}}1
      {ñ}{{\~n}}1 {Ñ}{{\~N}}1 {ü}{{\"u}}1 {Ü}{{\"U}}1
      {¿}{{?`}}1 {¡}{{!`}}1
      {—}{{---}}1 {–}{{--}}1 {−}{{-}}1
      {±}{{\ensuremath{\pm}}}1 {×}{{\ensuremath{\times}}}1
      {→}{{\ensuremath{\rightarrow}}}1
      {°}{{\textdegree}}1
      {·}{{\ensuremath{\cdot}}}1
      {€}{{\texteuro}}1
      {≈}{{\ensuremath{\approx}}}1
      {≤}{{\ensuremath{\le}}}1
      {⁰}{{\textsuperscript{0}}}1 {¹}{{\textsuperscript{1}}}1
      {²}{{\textsuperscript{2}}}1 {³}{{\textsuperscript{3}}}1
      {⁴}{{\textsuperscript{4}}}1 {⁵}{{\textsuperscript{5}}}1
      {⁶}{{\textsuperscript{6}}}1 {⁷}{{\textsuperscript{7}}}1
      {⁸}{{\textsuperscript{8}}}1 {⁹}{{\textsuperscript{9}}}1
}

\lstset{style=pythonstyle}

% Title slide macro: \lessontitle{Lesson Number}{Title}
\newcommand{\lessontitle}[2]{
    \title{Lesson #1: #2}
    \author{}
    \institute{Tec de Monterrey}
    \begin{frame}[plain]
        \titlepage
    \end{frame}
}

% Figure frame macro: \figureframe{Title}{Image path}
\newcommand{\figureframe}[2]{
    \begin{frame}{#1}
        \begin{center}
            \includegraphics[height=0.8\textheight,width=\textwidth,keepaspectratio]{#2}
        \end{center}
    \end{frame}
}


% Listings pulled from src/ with \lstinputlisting, so the slide is the file.
\lstset{
    basicstyle=\ttfamily\fontsize{8pt}{9.6pt}\selectfont,
    breaklines=true,
    breakatwhitespace=true,
    columns=fullflexible,
    keepspaces=true,
    xleftmargin=0.3em,
    framesep=2pt,
    aboveskip=0pt,
    belowskip=0pt
}

\newcommand{\resultframe}[3]{%
    \begin{frame}{#1}
        \begin{center}
            \includegraphics[height=0.62\textheight,width=\textwidth,keepaspectratio]{#2}
        \end{center}
        \vspace{0.25em}
        {\small #3\par}
    \end{frame}%
}

\newcommand{\claimframe}[2]{%
    \begin{frame}{#1}
        {\small #2\par}
    \end{frame}%
}

\date{}

\begin{document}

\lessontitle{09}{Binary Combinatorial}

\section{This lesson}

\begin{frame}{This lesson}
Lesson 08 mixed types in a box we could not draw. Today the chromosome is bits, and three problems share the operators.

\vspace{0.6em}
\begin{itemize}
    \item Knapsack, roster, radar: same bits, different meaning, different feasibility.
    \item Repair changes the distribution being searched; it is not free plumbing.
    \item Small instances should be checked exactly.
\end{itemize}
\end{frame}

% ================= knapsack_01_the_instance =================
\begin{frame}{Three binary models}
\small
\textbf{Knapsack:}\quad
\(\max\sum_i v_ix_i\) subject to \(\sum_i w_ix_i\le C\).\par
\medskip
\textbf{Radar:}\quad
\(\min\;U(x)(|S|+1)+\sum_i x_i\). Since at most \(|S|\) radars are active,
one fewer uncovered target dominates every radar-count change.\par
\medskip
\textbf{Scheduling:}\quad
\(\min\sum_{eds}c_{eds}x_{eds}\), with
\(\sum_e x_{eds}=d_s\) and \(\sum_{ds}x_{eds}\le5\).
Repair changes the distribution actually searched.
\end{frame}

\section{The instance}

\begin{frame}{The instance}
Each bit answers one question: take this item or leave it. Before introducing a
GA, enumerate this deliberately small instance so later runs have an honest
answer to compare against.
\end{frame}

\resultframe{The instance}{../figures/knapsack_01_the_instance.png}{%
    The 12-item instance has 4,096 chromosomes; exact optimum is value 50 at weight 20}

% ================= knapsack_02_random_population =================
\section{A random population}

\begin{frame}{A random population}
A random population
\end{frame}

\begin{frame}[fragile]{(1) Candidate}
\lstinputlisting[basicstyle=\ttfamily\fontsize{8pt}{9.6pt}\selectfont,firstline=59,lastline=69]{../src/knapsack_02_random_population.py}
\end{frame}

\begin{frame}[fragile]{(2) random\_candidate()}
\lstinputlisting[basicstyle=\ttfamily\fontsize{8pt}{9.6pt}\selectfont,firstline=72,lastline=77]{../src/knapsack_02_random_population.py}
\end{frame}

\begin{frame}[fragile]{(3) POPULATION\_SIZE}
\lstinputlisting[basicstyle=\ttfamily\fontsize{8pt}{9.6pt}\selectfont,firstline=80,lastline=82]{../src/knapsack_02_random_population.py}
\end{frame}

\begin{frame}[fragile]{(4) the census}
\lstinputlisting[basicstyle=\ttfamily\fontsize{8pt}{9.6pt}\selectfont,firstline=85,lastline=90]{../src/knapsack_02_random_population.py}
\end{frame}

\resultframe{A random population}{../figures/knapsack_02_random_population.png}{%
    233 of 1,000 random chromosomes are feasible; best random value is 47}

% ================= knapsack_03_repair =================
\section{Repair creates selection pressure}

\begin{frame}{Repair creates selection pressure}
Repair creates selection pressure
\end{frame}

\begin{frame}[fragile]{(1) repair()}
\lstinputlisting[basicstyle=\ttfamily\fontsize{8pt}{9.6pt}\selectfont,firstline=77,lastline=88]{../src/knapsack_03_repair.py}
\end{frame}

\begin{frame}[fragile]{(2) the comparison}
\lstinputlisting[basicstyle=\ttfamily\fontsize{8pt}{9.6pt}\selectfont,firstline=91,lastline=97]{../src/knapsack_03_repair.py}
\end{frame}

\resultframe{Repair creates selection pressure}{../figures/knapsack_03_repair.png}{%
    Repair makes 1,000/1,000 feasible and changes 767 chromosomes}

% ================= knapsack_04_selection_and_crossover =================
\section{Selection and crossover}

\begin{frame}{Selection and crossover}
Selection and crossover
\end{frame}

\begin{frame}[fragile]{(1) tournament()}
\lstinputlisting[basicstyle=\ttfamily\fontsize{8pt}{9.6pt}\selectfont,firstline=85,lastline=88]{../src/knapsack_04_selection_and_crossover.py}
\end{frame}

\begin{frame}[fragile]{(2) crossover()}
\lstinputlisting[basicstyle=\ttfamily\fontsize{8pt}{9.6pt}\selectfont,firstline=91,lastline=97]{../src/knapsack_04_selection_and_crossover.py}
\end{frame}

\begin{frame}[fragile]{(3) one\_generation()}
\lstinputlisting[basicstyle=\ttfamily\fontsize{8pt}{9.6pt}\selectfont,firstline=100,lastline=109]{../src/knapsack_04_selection_and_crossover.py}
\end{frame}

\resultframe{Selection and crossover}{../figures/knapsack_04_selection_and_crossover.png}{%
    One generation raises mean value from 39.00 to 40.73; all stored children remain feasible}

% ================= knapsack_05_the_full_search =================
\section{The full search}

\begin{frame}{The full search}
The full search
\end{frame}

\begin{frame}[fragile]{(1) mutate()}
\lstinputlisting[basicstyle=\ttfamily\fontsize{8pt}{9.6pt}\selectfont,firstline=97,lastline=102]{../src/knapsack_05_the_full_search.py}
\end{frame}

\begin{frame}[fragile]{(2) run()}
\lstinputlisting[basicstyle=\ttfamily\fontsize{8pt}{9.6pt}\selectfont,firstline=113,lastline=127]{../src/knapsack_05_the_full_search.py}
\end{frame}

\begin{frame}[fragile]{(3) the exact-gap report}
\lstinputlisting[basicstyle=\ttfamily\fontsize{8pt}{9.6pt}\selectfont,firstline=130,lastline=135]{../src/knapsack_05_the_full_search.py}
\end{frame}

\resultframe{The full search}{../figures/knapsack_05_the_full_search.png}{%
    The GA reaches the exact value 50 at weight 20 with zero gap in 4,060 evaluations}

% ================= radar_01_the_landscape =================
\section{A coverage landscape}

\begin{frame}{A coverage landscape}
Each bit decides whether to install a radar at one candidate site. A radar
covers a target when their Euclidean distance is at most the fixed radius.
\end{frame}

\resultframe{A coverage landscape}{../figures/radar_01_the_landscape.png}{%
    Nine sites define 512 plans and jointly cover all 36 targets}

% ================= radar_02_the_chromosome =================
\section{The chromosome}

\begin{frame}{The chromosome}
The chromosome
\end{frame}

\begin{frame}[fragile]{(1) Candidate}
\lstinputlisting[basicstyle=\ttfamily\fontsize{8pt}{9.6pt}\selectfont,firstline=39,lastline=49]{../src/radar_02_the_chromosome.py}
\end{frame}

\begin{frame}[fragile]{(2) random\_candidate()}
\lstinputlisting[basicstyle=\ttfamily\fontsize{8pt}{9.6pt}\selectfont,firstline=52,lastline=59]{../src/radar_02_the_chromosome.py}
\end{frame}

\begin{frame}[fragile]{(3) the census}
\lstinputlisting[basicstyle=\ttfamily\fontsize{8pt}{9.6pt}\selectfont,firstline=62,lastline=66]{../src/radar_02_the_chromosome.py}
\end{frame}

\resultframe{The chromosome}{../figures/radar_02_the_chromosome.png}{%
    21 of 500 random plans cover every target}

% ================= radar_03_the_penalty =================
\section{A lexicographic penalty}

\begin{frame}{A lexicographic penalty}
A lexicographic penalty
\end{frame}

\begin{frame}[fragile]{(1) objective()}
\lstinputlisting[basicstyle=\ttfamily\fontsize{8pt}{9.6pt}\selectfont,firstline=57,lastline=61]{../src/radar_03_the_penalty.py}
\end{frame}

\begin{frame}[fragile]{(2) the ranking audit}
\lstinputlisting[basicstyle=\ttfamily\fontsize{8pt}{9.6pt}\selectfont,firstline=64,lastline=70]{../src/radar_03_the_penalty.py}
\end{frame}

\resultframe{A lexicographic penalty}{../figures/radar_03_the_penalty.png}{%
    The audited penalty ranks every feasible plan ahead of every infeasible plan}

% ================= radar_04_operators =================
\section{Binary operators}

\begin{frame}{Binary operators}
Binary operators
\end{frame}

\begin{frame}[fragile]{(1) tournament()}
\lstinputlisting[basicstyle=\ttfamily\fontsize{8pt}{9.6pt}\selectfont,firstline=58,lastline=61]{../src/radar_04_operators.py}
\end{frame}

\begin{frame}[fragile]{(2) crossover()}
\lstinputlisting[basicstyle=\ttfamily\fontsize{8pt}{9.6pt}\selectfont,firstline=64,lastline=68]{../src/radar_04_operators.py}
\end{frame}

\begin{frame}[fragile]{(3) mutate()}
\lstinputlisting[basicstyle=\ttfamily\fontsize{8pt}{9.6pt}\selectfont,firstline=71,lastline=77]{../src/radar_04_operators.py}
\end{frame}

\resultframe{Binary operators}{../figures/radar_04_operators.png}{%
    One generation lowers mean objective from 86.0 to 42.6}

% ================= radar_05_the_full_search =================
\section{The full search}

\begin{frame}{The full search}
The full search
\end{frame}

\begin{frame}[fragile]{(1) exact\_optimum()}
\lstinputlisting[basicstyle=\ttfamily\fontsize{8pt}{9.6pt}\selectfont,firstline=76,lastline=80]{../src/radar_05_the_full_search.py}
\end{frame}

\begin{frame}[fragile]{(2) run()}
\lstinputlisting[basicstyle=\ttfamily\fontsize{8pt}{9.6pt}\selectfont,firstline=83,lastline=98]{../src/radar_05_the_full_search.py}
\end{frame}

\begin{frame}[fragile]{(3) the coverage map}
\lstinputlisting[basicstyle=\ttfamily\fontsize{8pt}{9.6pt}\selectfont,firstline=110,lastline=126]{../src/radar_05_the_full_search.py}
\end{frame}

\resultframe{The full search}{../figures/radar_05_the_full_search.png}{%
    GA and exhaustive enumeration both find five radars with full coverage}

% ================= schedule_01_the_requirements =================
\section{The requirements}

\begin{frame}{The requirements}
One bit assigns one employee to one shift. The instance needs 35 assignments
across 21 shifts, while eight employees can cover at most five shifts each.
\end{frame}

\resultframe{The requirements}{../figures/schedule_01_the_requirements.png}{%
    168 bits encode 35 required assignments under capacity 40}

% ================= schedule_02_random_rosters =================
\section{Random rosters}

\begin{frame}{Random rosters}
Random rosters
\end{frame}

\begin{frame}[fragile]{(1) random\_roster()}
\lstinputlisting[basicstyle=\ttfamily\fontsize{8pt}{9.6pt}\selectfont,firstline=34,lastline=39]{../src/schedule_02_random_rosters.py}
\end{frame}

\begin{frame}[fragile]{(2) violations()}
\lstinputlisting[basicstyle=\ttfamily\fontsize{8pt}{9.6pt}\selectfont,firstline=42,lastline=56]{../src/schedule_02_random_rosters.py}
\end{frame}

\begin{frame}[fragile]{(3) the census}
\lstinputlisting[basicstyle=\ttfamily\fontsize{8pt}{9.6pt}\selectfont,firstline=59,lastline=64]{../src/schedule_02_random_rosters.py}
\end{frame}

\resultframe{Random rosters}{../figures/schedule_02_random_rosters.png}{%
    0 of 500 random rosters are feasible; the best still has 14 violations}

% ================= schedule_03_repair =================
\section{Repair the constraints}

\begin{frame}{Repair the constraints}
Repair the constraints
\end{frame}

\begin{frame}[fragile]{(1) preference\_cost()}
\lstinputlisting[basicstyle=\ttfamily\fontsize{8pt}{9.6pt}\selectfont,firstline=55,lastline=59]{../src/schedule_03_repair.py}
\end{frame}

\begin{frame}[fragile]{(2) repair()}
\lstinputlisting[basicstyle=\ttfamily\fontsize{6pt}{7.2pt}\selectfont,firstline=62,lastline=94]{../src/schedule_03_repair.py}
\end{frame}

\begin{frame}[fragile]{(3) the comparison}
\lstinputlisting[basicstyle=\ttfamily\fontsize{8pt}{9.6pt}\selectfont,firstline=97,lastline=103]{../src/schedule_03_repair.py}
\end{frame}

\resultframe{Repair the constraints}{../figures/schedule_03_repair.png}{%
    Repair makes 500/500 rosters feasible; repaired costs range from 58 upward}

% ================= schedule_04_operators =================
\section{Operators under repair}

\begin{frame}{Operators under repair}
Operators under repair
\end{frame}

\begin{frame}[fragile]{(1) tournament()}
\lstinputlisting[basicstyle=\ttfamily\fontsize{8pt}{9.6pt}\selectfont,firstline=81,lastline=84]{../src/schedule_04_operators.py}
\end{frame}

\begin{frame}[fragile]{(2) crossover()}
\lstinputlisting[basicstyle=\ttfamily\fontsize{8pt}{9.6pt}\selectfont,firstline=87,lastline=92]{../src/schedule_04_operators.py}
\end{frame}

\begin{frame}[fragile]{(3) mutate()}
\lstinputlisting[basicstyle=\ttfamily\fontsize{8pt}{9.6pt}\selectfont,firstline=95,lastline=106]{../src/schedule_04_operators.py}
\end{frame}

\resultframe{Operators under repair}{../figures/schedule_04_operators.png}{%
    All 100 children are feasible; mean cost falls from 77.8 to 61.2}

% ================= schedule_05_the_full_search =================
\section{The full search}

\begin{frame}{The full search}
The full search
\end{frame}

\begin{frame}[fragile]{(1) one\_generation()}
\lstinputlisting[basicstyle=\ttfamily\fontsize{8pt}{9.6pt}\selectfont,firstline=102,lastline=109]{../src/schedule_05_the_full_search.py}
\end{frame}

\begin{frame}[fragile]{(2) run()}
\lstinputlisting[basicstyle=\ttfamily\fontsize{8pt}{9.6pt}\selectfont,firstline=112,lastline=123]{../src/schedule_05_the_full_search.py}
\end{frame}

\begin{frame}[fragile]{(3) the baseline report}
\lstinputlisting[basicstyle=\ttfamily\fontsize{8pt}{9.6pt}\selectfont,firstline=126,lastline=131]{../src/schedule_05_the_full_search.py}
\end{frame}

\resultframe{The full search}{../figures/schedule_05_the_full_search.png}{%
    GA cost 41 beats the best of 5,100 repaired-random rosters, cost 52, at equal evaluation count}

\section{Next}

\begin{frame}{Next}
The knapsack GA hits the exact value 50 at weight 20. The roster starts at 0/500 feasible and still beats repaired-random at equal budget.

\vspace{0.8em}
\textbf{Lesson 10 --- Ordered Combinatorial Optimisation}

Order and uniqueness are now part of correctness. The OX1 that was a silent no-op in Lesson 04 is the legal operator here.
\end{frame}
\end{document}
